% ==========================================================================
%  Chapter 65 — Ko–Bathe Concrete Model: 3D Extension (N=6)
% ==========================================================================

\chapter{Ko--Bathe Concrete Model: Three-Dimensional Extension}
\label{ch:ko-bathe-3d}

Chapter~\ref{ch:ko-bathe-concrete} documented the plane-stress ($N\!=\!3$)
implementation of the Ko--Bathe plastic-fracturing concrete model.  This
chapter extends the formulation to the full three-dimensional stress space
($N\!=\!6$), following the same algorithmic structure---fracturing, plasticity,
cracking---but with the geometric and algebraic modifications required by six
independent stress components and three potential crack planes.

The implementation resides in a single header,
\texttt{KoBatheConcrete3D.hh} ($\sim$620~lines), and reuses the \texttt{KoBatheParameters}
structure from the 2D version without modification.

% ========================================================================
\section{Voigt Convention and State Variables}
\label{sec:kb3d-voigt}

The Voigt ordering follows the \texttt{fall\_n} convention:
\begin{equation}
  \bsigma = (\sigma_{xx},\;\sigma_{yy},\;\sigma_{zz},\;
              \tau_{yz},\;\tau_{xz},\;\tau_{xy})^T, \quad
  \bvarepsilon = (\varepsilon_{xx},\;\varepsilon_{yy},\;\varepsilon_{zz},\;
                   \gamma_{yz},\;\gamma_{xz},\;\gamma_{xy})^T,
  \label{eq:voigt-3d}
\end{equation}
where $\gamma_{ij} = 2\varepsilon_{ij}$ denotes engineering shear strain.

The internal state structure \texttt{KoBatheState3D} stores:
\begin{itemize}
  \item \textbf{Fracture history}: $\sigma_{o,\max}$ and $\tau_{o,\max}$
        (octahedral stress maxima driving stiffness degradation).
  \item \textbf{Plastic strain}: $\bvarepsilon^p \in \mathbb{R}^6$ in Voigt
        notation, plus three effective plastic strains $\bar\varepsilon^p_1$,
        $\bar\varepsilon^p_2$, $\bar\varepsilon^p_3$ (compression coordinate
        hardening variables).
  \item \textbf{Crack data}: up to 3 orthogonal smeared cracks, each defined
        by a normal vector $\mathbf{n}_i \in \mathbb{R}^3$, opening strain
        $\varepsilon_{nn,i}$, maximum historical opening
        $\varepsilon_{nn,i}^{\max}$, and a closure flag.
\end{itemize}


% ========================================================================
\section{Octahedral Stress Invariants in 3D}
\label{sec:kb3d-octahedral}

For a general three-dimensional stress state with principal stresses
$\sigma_1 \ge \sigma_2 \ge \sigma_3$ obtained from the eigenvalue
decomposition of the Cauchy stress tensor, the octahedral invariants are

\begin{align}
  \sigma_o &= \frac{1}{3}(\sigma_1 + \sigma_2 + \sigma_3), \label{eq:sigma-o-3d} \\[4pt]
  \tau_o   &= \frac{\sqrt{2\,J_2}}{3}, \qquad
  J_2 = \tfrac{1}{2}(s_1^2 + s_2^2 + s_3^2), \label{eq:tau-o-3d}
\end{align}
where $s_i = \sigma_i - \sigma_o$ are deviatoric principal stresses.

The Lode angle $\theta \in [0, \pi/3]$ is computed via
\begin{equation}
  \cos(3\theta) = \frac{3\sqrt{3}}{2}\,\frac{J_3}{J_2^{3/2}},
  \qquad J_3 = s_1\,s_2\,s_3.
  \label{eq:lode-angle}
\end{equation}
The Lode angle governs the shape of the failure surface in the deviatoric
plane through an elliptic interpolation function (see
Section~\ref{sec:kb3d-cracking}).

\paragraph{Reduction to plane stress.}
Setting $\sigma_3 = 0$ and computing the principal stresses from
$(\sigma_{xx}, \sigma_{yy}, \tau_{xy})$ recovers the 2D octahedral invariants
of Chapter~\ref{ch:ko-bathe-concrete}.


% ========================================================================
\section{Principal Strain Decomposition}
\label{sec:kb3d-principal}

The elastic strain tensor
\[
  \mathbf{E}^e =
  \begin{pmatrix}
    \varepsilon_{xx} & \varepsilon_{xy} & \varepsilon_{xz} \\
    \varepsilon_{xy} & \varepsilon_{yy} & \varepsilon_{yz} \\
    \varepsilon_{xz} & \varepsilon_{yz} & \varepsilon_{zz}
  \end{pmatrix},
  \qquad \varepsilon_{ij} = \gamma_{ij}/2 \text{ for } i \neq j,
\]
is decomposed by Eigen's \texttt{SelfAdjointEigenSolver} into sorted principal
values $\varepsilon_1 \ge \varepsilon_2 \ge \varepsilon_3$ and an orthogonal
eigenvector matrix~$\mathbf{V}$.  The rows of $\mathbf{V}$ are the principal
directions, so that the rotation from global to principal coordinates is
$\varepsilon^{\text{princ}} = \mathbf{V} \cdot \varepsilon^{\text{global}}
 \cdot \mathbf{V}^T$.


% ========================================================================
\section{Compression Coordinates}
\label{sec:kb3d-compression}

Plasticity operates in a \emph{compression coordinate} space derived from the
compressive magnitudes of the principal elastic strains:
\begin{equation}
  D_i = \max(-\varepsilon_i^e, \; 0), \qquad i = 1, 2, 3.
  \label{eq:Di}
\end{equation}
Since $\varepsilon_1 \ge \varepsilon_2 \ge \varepsilon_3$, we have
$D_3 \ge D_2 \ge D_1 \ge 0$.

The three elastic strain coordinates that drive the yield surfaces are:
\begin{align}
  \hat{e}_1 &= \frac{D_1 + D_2 + D_3}{\sqrt{3}}
             && \text{(hydrostatic)}, \label{eq:ee1-3d} \\
  \hat{e}_2 &= D_3 - D_1
             && \text{(deviatoric)}, \label{eq:ee2-3d} \\
  \hat{e}_3 &= \frac{D_1 + D_2 + D_3}{2}
             && \text{(biaxial)}. \label{eq:ee3-3d}
\end{align}

\paragraph{Consistency with the 2D formulation.}
In the plane-stress case ($\varepsilon_3 = \sigma_{zz} = 0$, so
$D_3 = 0$ whenever $\varepsilon_3 \ge 0$), the coordinates reduce to
\[
  \hat{e}_1 \to \frac{D_1 + D_2}{\sqrt{3}}, \qquad
  \hat{e}_2 \to D_2 - D_1, \qquad
  \hat{e}_3 \to \frac{D_1 + D_2}{2},
\]
which is exactly the 2D definition in Chapter~\ref{ch:ko-bathe-concrete}
(where $\tilde{e}_1 = D_2 - D_1$ and $\tilde{e}_2 = D_1 + D_2$, mapped to
$\hat{e}_1 = \tilde{e}_2/\sqrt{3}$, $\hat{e}_2 = \tilde{e}_1$,
$\hat{e}_3 = \tilde{e}_2/2$).

This definition—which uses the simple difference $D_3 - D_1$ rather than an
invariant-based deviatoric radius $\sqrt{2J_{2,D}}$---ensures that the
calibration constants ($c_{p2} = 36$, $c_{p3} = 12$, etc.)\ obtained for
the 2D model remain valid in 3D without re-calibration.


% ========================================================================
\section{Isotropic 3D Tangent Matrix}
\label{sec:kb3d-tangent}

The secant and tangent constitutive matrices use the standard isotropic form
\begin{equation}
  \mathbf{C}(K, G) =
  \begin{pmatrix}
    \lambda + 2\mu & \lambda & \lambda & 0 & 0 & 0 \\
    \lambda & \lambda + 2\mu & \lambda & 0 & 0 & 0 \\
    \lambda & \lambda & \lambda + 2\mu & 0 & 0 & 0 \\
    0 & 0 & 0 & G & 0 & 0 \\
    0 & 0 & 0 & 0 & G & 0 \\
    0 & 0 & 0 & 0 & 0 & G
  \end{pmatrix}, \quad
  \lambda = K - \tfrac{2}{3}G, \quad
  \mu = G.
  \label{eq:C-iso-3d}
\end{equation}
Here $K$ and $G$ are the \emph{secant} bulk and shear moduli from the
fracturing algorithm (identical scalar formulas as in 2D) or the
\emph{consistent} (tangent) moduli $K_c$, $G_c$.


% ========================================================================
\section{Mandel Rotation for the Cracked Tangent}
\label{sec:kb3d-mandel}

In 2D, applying crack modifications to the tangent matrix requires only a
$3\times3$ rotation.  In 3D, the analogous transformation involves a
$6\times6$ matrix whose construction requires careful handling of the
engineering/tensor shear distinction.

\subsection{Mandel Notation}

The \emph{Mandel} (or normalised Voigt) representation scales the shear
components by $\sqrt{2}$:
\[
  \boldsymbol{\sigma}^M =
  (\sigma_{11},\;\sigma_{22},\;\sigma_{33},\;
   \sqrt{2}\,\sigma_{23},\;\sqrt{2}\,\sigma_{13},\;\sqrt{2}\,\sigma_{12})^T.
\]
In this notation, a rotation $\mathbf{R} \in \mathrm{SO}(3)$ applied to a
second-order tensor corresponds to an \emph{orthogonal} $6\times6$ matrix
$\mathbf{Q}$ acting on the Mandel vector:
\[
  \boldsymbol{\sigma}^M_{\text{rot}} = \mathbf{Q}\,\boldsymbol{\sigma}^M, \qquad
  \mathbf{Q}^{-1} = \mathbf{Q}^T.
\]
The orthogonality $\mathbf{Q}^T\mathbf{Q} = \mathbf{I}$ is the key property
that makes Mandel notation attractive: fourth-order tensor transformations
reduce to simple matrix products.

\subsection{Construction of the Mandel Rotation Matrix}

Given a rotation matrix $\mathbf{R}$ whose rows are the basis vectors
$(\mathbf{n}, \mathbf{t}_1, \mathbf{t}_2)$ of the local crack frame, the
$6\times6$ Mandel rotation matrix has four blocks:

\paragraph{Normal--Normal block} ($3\times3$, upper-left):
\[
  Q_{ij} = R_{ij}^2, \qquad i, j \in \{0, 1, 2\}.
\]

\paragraph{Normal--Shear block} ($3\times3$, upper-right):
For shear index $s \leftrightarrow (a, b)$ with convention
$3 \to (1,2)$, $4 \to (0,2)$, $5 \to (0,1)$:
\[
  Q_{i,\,s+3} = \sqrt{2}\, R_{i,a}\, R_{i,b}.
\]

\paragraph{Shear--Normal block} ($3\times3$, lower-left):
\[
  Q_{s+3,\,j} = \sqrt{2}\, R_{a,j}\, R_{b,j}.
\]

\paragraph{Shear--Shear block} ($3\times3$, lower-right):
\[
  Q_{s_1+3,\,s_2+3} = R_{a_1, a_2}\, R_{b_1, b_2} + R_{a_1, b_2}\, R_{b_1, a_2}.
\]

\subsection{Engineering--Mandel Conversion}

Conversion between the engineering Voigt matrix $\mathbf{C}^{\text{eng}}$ and
the Mandel matrix $\mathbf{C}^M$ is performed via a diagonal scaling matrix
$\mathbf{D} = \mathrm{diag}(1, 1, 1, \sqrt{2}, \sqrt{2}, \sqrt{2})$:
\begin{equation}
  \mathbf{C}^M = \mathbf{D}\,\mathbf{C}^{\text{eng}}\,\mathbf{D}^{-1}.
  \label{eq:eng-mandel}
\end{equation}

\subsection{Cracked Tangent Application}

For each active crack $i$ with normal $\mathbf{n}_i$:
\begin{enumerate}
  \item Build an orthonormal basis $(\mathbf{n}, \mathbf{t}_1, \mathbf{t}_2)$.
  \item Construct the rotation matrix $\mathbf{R}$ with rows
        $(\mathbf{n}, \mathbf{t}_1, \mathbf{t}_2)$.
  \item Convert $\mathbf{C}^{\text{eng}}$ to Mandel:
        $\mathbf{C}^M = \mathbf{D}\mathbf{C}\mathbf{D}^{-1}$.
  \item Rotate to crack-local frame:
        $\mathbf{C}^M_{\text{loc}} = \mathbf{Q}\,\mathbf{C}^M\,\mathbf{Q}^T$.
  \item In the local frame, modify the $(0,0)$ component (normal stiffness):
        \[
          C^M_{\text{loc}}(0,0) \leftarrow
          \begin{cases}
            C_{\text{intact}} & \text{if crack closed}, \\
            C_{ts}           & \text{if in softening zone } (\varepsilon_{tp} \le \varepsilon_{nn} \le \varepsilon_{tu}), \\
            \eta_N \cdot C_{\text{intact}} & \text{otherwise (fully open)}.
          \end{cases}
        \]
        Decouple the normal row/column: $C^M(0,j) = C^M(j,0) = 0$ for $j \neq 0$.
        Apply shear retention: $C^M(4,4) \leftarrow \eta_S \cdot C^M(4,4)$,
        $C^M(5,5) \leftarrow \eta_S \cdot C^M(5,5)$.
  \item Rotate back: $\mathbf{C}^M = \mathbf{Q}^T\,\mathbf{C}^M_{\text{loc}}\,\mathbf{Q}$.
  \item Convert back to engineering: $\mathbf{C}^{\text{eng}} = \mathbf{D}^{-1}\mathbf{C}^M\mathbf{D}$.
\end{enumerate}

Multiple cracks are applied sequentially; each crack modifies the tangent
produced by the previous one.


% ========================================================================
\section{Crack Initiation in 3D}
\label{sec:kb3d-cracking}

A new crack is initiated when the maximum principal stress direction exceeds
the crack function envelope.  The crack function uses the Menetr\'{e}y--Willam
form with the Lode angle $\theta$ from Eq.~\eqref{eq:lode-angle}:

\begin{equation}
  g(\sigma_o, \tau_o, \theta) = \tau_o - \frac{0.944\,f_c'}{w(e, \theta)} \,
  \left(t_p - \frac{\sigma_o}{f_c'}\right)^{0.724},
  \label{eq:crack-function-3d}
\end{equation}
where $w(e, \theta)$ is the polar radius of an elliptic deviatoric cross-section:
\[
  w(e,\theta) = \frac{4(1-e^2)\cos^2\theta + (2e-1)^2}
  {2(1-e^2)\cos\theta + (2e-1)\sqrt{4(1-e^2)\cos^2\theta + 5e^2 - 4e}},
\]
and the eccentricity $e$ depends on the current stress state:
\begin{equation}
  e = 0.670551 \left(t_p - \frac{\sigma_o}{f_c'}\right)^{0.133}.
\end{equation}

A crack initiates when $g > 0$.  Its normal is the eigenvector corresponding
to the largest principal stress.  Additional cracks are orthogonalised:
\begin{itemize}
  \item The second crack normal is projected to lie in the plane
        perpendicular to the first crack normal.
  \item The third crack normal is computed as the cross product of the
        first two normals.
\end{itemize}
This guarantees a mutually orthogonal set of up to three crack planes in 3D.


% ========================================================================
\section{Plasticity Return Mapping}
\label{sec:kb3d-plasticity}

The plasticity algorithm operates on the same three yield surfaces as the 2D
model (hydrostatic, uniaxial, biaxial), defined in the compression coordinate
space of Section~\ref{sec:kb3d-compression}:
\begin{align}
  f_1 &= \hat{e}_1 - f_c'\,f_{p1}(\bar\varepsilon^p_1), \label{eq:f1-3d} \\
  f_2 &= \hat{e}_2 - f_c'\,f_{p2}(\bar\varepsilon^p_2), \\
  f_3 &= \hat{e}_3 - f_c'\,f_{p3}(\bar\varepsilon^p_3). \label{eq:f3-3d}
\end{align}
The failure curves $f_{p1}$, $f_{p2}$, $f_{p3}$ and their derivatives are
identical to those in 2D (Chapter~\ref{ch:ko-bathe-concrete},
Section~``Yield Functions'').

\subsection{Back-Transform: Compression Coordinates to Global Voigt}

The Newton iteration produces plastic multiplier increments
$\Delta\lambda_1$ (hydrostatic), $\Delta\lambda_2$ (deviatoric), and
$\Delta\lambda_3$ (biaxial; only updates the hardening variable
$\bar\varepsilon^p_3$, no plastic flow).

The back-transform from compression coordinates to principal plastic strain
increments is:
\begin{align}
  \Delta\varepsilon^p_1 &= -\!\left(\frac{\Delta\lambda_1}{\sqrt{3}} - \frac{\Delta\lambda_2}{2}\right), \\
  \Delta\varepsilon^p_2 &= -\frac{\Delta\lambda_1}{\sqrt{3}}, \\
  \Delta\varepsilon^p_3 &= -\!\left(\frac{\Delta\lambda_1}{\sqrt{3}} + \frac{\Delta\lambda_2}{2}\right).
\end{align}
This is the natural 3D extension of the 2D back-transform: the volumetric
contribution $\Delta\lambda_1/\sqrt{3}$ is shared equally among all three
principal directions, while the deviatoric contribution $\pm\Delta\lambda_2/2$
is added to the extreme directions ($D_1$ and $D_3$) symmetrically.

\paragraph{Principal to global.}
The principal plastic strain increment tensor
$\Delta\varepsilon^p_{\text{princ}} = \mathrm{diag}(\Delta\varepsilon^p_1,
\Delta\varepsilon^p_2, \Delta\varepsilon^p_3)$
is rotated to the global frame using the eigenvector matrix $\mathbf{V}$:
\[
  \Delta\boldsymbol{\varepsilon}^p_{\text{global}} = \mathbf{V}^T \cdot
  \Delta\boldsymbol{\varepsilon}^p_{\text{princ}} \cdot \mathbf{V},
\]
and the resulting tensor is mapped to Voigt form (with engineering shear
$\gamma_{ij} = 2\varepsilon_{ij}$).


% ========================================================================
\section{Full Algorithm}
\label{sec:kb3d-algorithm}

The 3D evaluation follows the same eight-step structure as the 2D model
(Table~1 in~\cite{KoBathe2026}), with the dimensional extensions described
above:

\begin{enumerate}
  \item Elastic strain: $\bvarepsilon^e = \bvarepsilon - \bvarepsilon^p$.
  \item Principal decomposition of $\bvarepsilon^e$ (eigenvalues + eigenvectors).
  \item Fracture moduli $K_s$, $G_s$, $\sigma_{id}$ from history
        ($\sigma_{o,\max}$, $\tau_{o,\max}$).
  \item Secant stress: $\bsigma = \mathbf{C}(K_s, G_s)\,\bvarepsilon^e
        - \sigma_{id}\,\boldsymbol{\delta}$.
  \item Update fracture history (track envelope of $\sigma_o$, $\tau_o$).
  \item Cracking check: evaluate crack function~\eqref{eq:crack-function-3d};
        if $g > 0$ and fewer than 3 cracks exist, initiate new crack with
        orthogonality enforcement.
  \item Update crack opening strains $\varepsilon_{nn,i}$; apply cracked
        tangent (Section~\ref{sec:kb3d-mandel}); recompute stress.
  \item Plasticity return mapping
        (Section~\ref{sec:kb3d-plasticity}); update $\bvarepsilon^p$ and
        recalculate equilibrium stress.
\end{enumerate}


% ========================================================================
\section{Integration with the Sub-Model Solver}
\label{sec:kb3d-integration}

The 3D model is integrated into the multiscale sub-model solver chain
through the \texttt{fall\_n} type-erasure infrastructure:

\begin{lstlisting}
InelasticMaterial<KoBatheConcrete3D> mat_inst{fc_MPa};
Material<ThreeDimensionalMaterial> mat{mat_inst, InelasticUpdate{}};
Model<ThreeDimensionalMaterial, SmallStrain, 3> M{domain, mat};
\end{lstlisting}

Here \texttt{InelasticMaterial<R>} is a type alias for
\texttt{MaterialInstance<R, CommittedState>}, which wraps the constitutive
relation \texttt{R} in an algorithmic state holder that provides the
\texttt{commit()} / \texttt{revert()} interface required by the Newton--Raphson
solver.

\subsection{Operator-Splitting Crack Check}
\label{sec:kb3d-operator-splitting}

A critical convergence issue arises from the sharp transition between
\emph{uncracked} ($\mathbf{C} = \mathbf{C}_e$) and \emph{cracked}
($\mathbf{C} = \eta_N\,\mathbf{C}_e$) tangent stiffness.  When several
Gauss points are near the crack-initiation threshold, the standard
Newton procedure oscillates: an element cracks $\to$ stress drops $\to$
Newton update reverses the strain $\to$ element un-cracks $\to$ stress
rises $\to$ element cracks again.  This creates a non-convergent 2-cycle
that no amount of damping, line-search tuning, or bisection can resolve.

The \texttt{evaluate()} function is pure-functional (copies the committed
state and returns a modified copy without mutating the input), so the
oscillation is between consistent but \emph{different} branches of the
stress function at each iteration---not from state mutation.

The fix is an \emph{operator splitting} strategy: the
\texttt{check\_new\_cracks} flag (default \texttt{true}) in
\texttt{evaluate()} gates the crack-initiation check (Step~7).
During Newton iterations, \texttt{compute\_response()} and
\texttt{tangent()} pass \texttt{false}; only \texttt{commit()} passes
\texttt{true}.  This means:

\begin{itemize}
  \item During Newton iterations: the stress function $\bsigma(\bvarepsilon)$
        is smooth---it uses only the \emph{committed} crack state (no new
        cracks form).  Newton converges reliably.
  \item At commit time: \texttt{evaluate()} runs with full crack checking,
        propagating any new cracks into the committed state.
  \item The next load step begins from the updated state, and Newton
        iteration proceeds smoothly again.
\end{itemize}

\subsection{PETSc Solver Configuration}

The sub-model solver uses incremental displacement control with 20 steps
and 8 levels of recursive bisection:

\begin{lstlisting}
PetscOptionsSetValue(nullptr, "-snes_linesearch_type", "basic");
PetscOptionsSetValue(nullptr, "-snes_max_it", "50");
PetscOptionsSetValue(nullptr, "-ksp_type", "preonly");
PetscOptionsSetValue(nullptr, "-pc_type", "lu");
constexpr int NUM_STEPS  = 20;
constexpr int MAX_BISECT = 8;
\end{lstlisting}

Key design choices:

\begin{itemize}
  \item \textbf{Basic line search} (no backtracking): the committed-crack
        tangent still creates piecewise-smooth residuals that confuse
        the backtracking line search (PETSc reason $-6$ =
        \texttt{SNES\_DIVERGED\_LINE\_SEARCH} in v3.24).  Basic line search
        takes the full Newton step and relies on the operator-splitting
        smoothness for convergence.
  \item \textbf{Direct LU}: always succeeds for the sub-model problem
        size ($\sim$1350 DOFs).
  \item \textbf{20 steps with bisection}: each step imposes $\Delta p = 0.05$.
        At each commit boundary new cracks form, and the next step starts
        from a well-defined state.
\end{itemize}


% ========================================================================
\section{Test Suite}
\label{sec:kb3d-tests}

The test file \texttt{tests/test\_ko\_bathe\_concrete\_3d.cpp} validates 23
checks across 10 test groups:

\begin{enumerate}
  \item \textbf{Concept satisfaction} (3 compile-time \texttt{static\_assert}s).
  \item \textbf{3D elastic response}: tangent symmetry, positive diagonal,
        correct $\sigma_{xx}$, $\sigma_{yy}$ for uniaxial strain.
  \item \textbf{Uniaxial compression} curves for $f_c' = 21$ and $30$~MPa
        (CSV export); peak stress within calibration bounds.
  \item \textbf{Triaxial compression}: confinement increases peak mean stress
        beyond $f_c'$.
  \item \textbf{Tension cracking}: positive peak tensile stress $\le 2f_t$;
        crack normal aligned with the tension direction; 1 crack formed.
  \item \textbf{Mandel rotation orthogonality}: hydrostatic strain produces
        equal normal stresses and vanishing shear.
  \item \textbf{Biaxial compression}: enhanced strength beyond $0.5\,f_c'$.
  \item \textbf{Commit cycle consistency}: external-state and internal-state
        evaluation paths produce identical stresses.
  \item \textbf{Pure shear}: positive $\tau_{xy}$, normal stresses $\approx 0$.
  \item \textbf{3D vs.\ 2D consistency}: uniaxial stress path comparison; maximum
        stress deviation within $50\%$ of $f_c'$.
\end{enumerate}

All 23/23 checks pass.  The CSV output for the uniaxial compression curves
allows visual comparison with the 2D model and the reference
paper~\cite{KoBathe2026}.


% ========================================================================
\section{Discussion}
\label{sec:kb3d-discussion}

\paragraph{Compression coordinate choice.}
An early implementation used $\hat{e}_2 = \sqrt{2\,J_{2,D}}$ (the deviatoric
radius of the compressive strain vector $\mathbf{D}$) as a more
``invariant-based'' extension.  However, this yields
$\hat{e}_2 \approx 0.82\,\varepsilon_c$ under uniaxial compression (versus
$\hat{e}_2 = \varepsilon_c$ with the current definition), delaying plasticity
activation and producing peak stresses $\sim$\!3$\times f_c'$.  The direct
difference $D_3 - D_1$ preserves consistency with the original 2D calibration
constants.

\paragraph{Tangent conditioning and stiffness retention.}
The cracked tangent features stiffness retention factors
$\eta_N = 0.20$ (normal) and $\eta_S = 0.50$ (shear), significantly
larger than the 2D counterparts ($10^{-4}$ and $0.1$).  The higher
values improve SNES convergence without significantly affecting the
structural response, since the sub-model elements are small ($\sim$10\,cm)
and the global behaviour is governed by the beam kinematics at the
structural scale.

\paragraph{Newton convergence and operator splitting.}
The ``pure-functional'' design of \texttt{evaluate()} introduces a subtle
convergence pathology.  Because \texttt{compute\_response()} always starts
from the \emph{committed} state (before the current load step), each
Newton iteration re-evaluates the stress without memory of what happened
in the previous iteration.  Near the cracking threshold, this creates a
non-convergent 2-cycle: trial strain crosses the threshold $\to$ crack
forms (in the local copy) $\to$ stress drops $\to$ Newton update reverses
the strain $\to$ no crack forms $\to$ stress rises $\to$ cycle repeats.

This 2-cycle manifests as PETSc convergence reason $-6$
(\texttt{SNES\_DIVERGED\_LINE\_SEARCH} in v3.24, \emph{not}
\texttt{SNES\_DIVERGED\_LINEAR\_SOLVE}, which is reason $-3$).  The
backtracking line search fails because the residual oscillates rather
than decreasing monotonically.

The fix is operator splitting: the flag
\texttt{check\_new\_cracks\,=\,false} is passed during Newton iterations
(via \texttt{compute\_response} and \texttt{tangent}), while
\texttt{commit()} calls with \texttt{true} to allow new cracks to form
only after convergence.  Combined with the basic (full-step) line search,
typical sub-model solves converge in 3--34 SNES iterations per load
increment, with occasional bisections near $p \approx 0.95$.

\paragraph{Mandel vs.\ engineering Voigt.}
The Mandel notation transformation is central to the 3D cracked tangent.
Using engineering Voigt directly would require separate transformation
matrices for stress and strain ($\mathbf{T}_\sigma \neq \mathbf{T}_\varepsilon$),
and the product
$\mathbf{C}_{\text{rot}} = \mathbf{T}_\sigma^{-1}\,\mathbf{C}_{\text{loc}}\,\mathbf{T}_\varepsilon$
loses the symmetry guarantee.  Mandel notation avoids this by ensuring
$\mathbf{Q}^{-1} = \mathbf{Q}^T$, so that
$\mathbf{C}^M_{\text{rot}} = \mathbf{Q}^T\,\mathbf{C}^M_{\text{loc}}\,\mathbf{Q}$
is symmetric whenever $\mathbf{C}^M_{\text{loc}}$ is.
